\documentclass[a4paper,12pt]{article}

\usepackage[utf8]{inputenc}
\usepackage[T1]{fontenc}
\usepackage[turkish]{babel}
\usepackage{geometry}
\usepackage{graphicx}
\usepackage{lmodern}
\usepackage{booktabs}
\usepackage{array}
\usepackage{setspace}
\usepackage{parskip}

\geometry{
    a4paper,
    left=3cm,
    right=2.5cm,
    top=2cm,
    bottom=2cm,
    headheight=15pt
}

\setstretch{1.2}
\setlength{\parskip}{0.8em}
\setlength{\parindent}{0pt}

\begin{document}

% === KAYNAKÇA ===
\clearpage
\section*{KAYNAKÇA}
\addcontentsline{toc}{section}{KAYNAKÇA}

\vspace{1em}
\begin{itemize}
    \item[] [1] T. Koshy, \textit{Elementary Number Theory with Applications}, 2nd Edition, Elsevier Science, 2007.
\end{itemize}

% === ÖZGEÇMİŞ ===
\clearpage
\section*{ÖZGEÇMİŞ}
\addcontentsline{toc}{section}{ÖZGEÇMİŞ}

2002 yılında İstanbul’da doğdu. İlköğrenimini Beşiktaş İlköğretim Okulu’nda, 
ortaöğrenimini Şair Nedim Ortaokulu’nda tamamladı. Lise eğitimine Erhan Gedikbaşı Anadolu Lisesi’nde başlayıp, 
11. sınıftan itibaren Uğur Anadolu Lisesi’nde devam ederek mezun oldu.

2021 yılında İstanbul Kültür Üniversitesi Fen-Edebiyat Fakültesi Matematik ve Bilgisayar Bilimleri Bölümü’ne kaydoldu ve lisans eğitimine burada devam etmektedir.

Akademik olarak sayılar teorisi, matematiksel mantık ve algoritmalarla ilgilenmektedir. Lisans eğitimi süresince çeşitli matematiksel modelleme projelerinde yer almış ve sayılar teorisi üzerine yoğunlaşan bir bitirme projesi hazırlamıştır.

İyi derecede İngilizce bilmektedir. Python, C/C++ ve LaTeX gibi programlama ve dokümantasyon araçlarında bilgi sahibidir.



\end{document}
